\documentclass[10pt,letterpaper,oneside]{report}
\usepackage[T1]{fontenc}
\usepackage[utf8]{inputenc}
\usepackage{lmodern}
\usepackage{microtype}
\usepackage[margin=0.85in,headheight=15pt]{geometry}
\usepackage{enumitem,booktabs,longtable,tabularx,array,makecell}
\usepackage{fancyhdr,lastpage,xcolor,hyperref,titlesec}
\usepackage{listings}
\usepackage[most]{tcolorbox}
\definecolor{ink}{HTML}{182D40}
\definecolor{accent}{HTML}{126B77}
\definecolor{pale}{HTML}{EEF5F5}
\definecolor{grayblue}{HTML}{677887}
\definecolor{codebg}{HTML}{F5F7F8}
\hypersetup{colorlinks=true,linkcolor=accent,urlcolor=accent,pdfauthor={C23 Engineering Companion},pdfsubject={ISO/IEC 9899:2024 educational engineering companion}}
\titleformat{\chapter}[display]{\normalfont\huge\bfseries\color{ink}}{\filleft\large\color{accent}CHAPTER \thechapter}{1ex}{\titlerule[1pt]\vspace{1ex}}
\titleformat{\section}{\Large\bfseries\color{ink}}{\thesection}{0.7em}{}
\titleformat{\subsection}{\large\bfseries\color{accent}}{\thesubsection}{0.65em}{}
\setlist[itemize]{leftmargin=1.3em,itemsep=2pt,topsep=3pt}
\setlist[enumerate]{leftmargin=1.6em,itemsep=3pt,topsep=3pt}
\setlength{\parskip}{5pt}
\setlength{\parindent}{0pt}
\renewcommand{\arraystretch}{1.18}
\pagestyle{fancy}\fancyhf{}\fancyhead[L]{\small\sffamily C23 ENGINEERING SERIES}\fancyhead[R]{\small\sffamily\color{accent}\DocShort}\fancyfoot[L]{\small\color{grayblue}ISO/IEC 9899:2024 companion}\fancyfoot[R]{\small\thepage\ / \pageref{LastPage}}
\newtcolorbox{keybox}{breakable,colback=pale,colframe=accent,boxrule=.6pt,arc=1mm,left=8pt,right=8pt,top=6pt,bottom=6pt}
\newtcolorbox{warnbox}{breakable,colback=codebg,colframe=grayblue,boxrule=.6pt,arc=1mm,left=8pt,right=8pt,top=6pt,bottom=6pt}
\lstdefinestyle{cstyle}{language=C,basicstyle=\ttfamily\footnotesize,keywordstyle=\color{accent}\bfseries,commentstyle=\color{grayblue}\itshape,stringstyle=\color{ink},backgroundcolor=\color{codebg},frame=single,rulecolor=\color{grayblue!40},breaklines=true,showstringspaces=false,tabsize=4,columns=fullflexible,keepspaces=true,upquote=true}
\lstset{style=cstyle}
\newcommand{\refcl}[1]{\textsf{\footnotesize[ISO/IEC 9899:2024, #1]}}
\newcommand{\standard}{\textbf{[STANDARD]}}
\newcommand{\optional}{\textbf{[OPTIONAL]}}
\newcommand{\guidance}{\textbf{[GUIDANCE]}}
\newcommand{\testtag}{\textbf{[TEST]}}
\newcommand{\term}[1]{\textbf{#1}}
\newcommand{\DocShort}{SECURE C HANDBOOK}
\begin{document}
\begin{titlepage}
\color{ink}\vspace*{1.2in}
{\sffamily\small\color{accent}\bfseries TECHNICAL DOCUMENTATION SUITE \hfill EDITION 2026-10}\par
\vspace{.3in}\rule{\linewidth}{1.5pt}\par\vspace{.22in}
{\Huge\bfseries Secure C23 Programming Handbook}\par\vspace{.22in}
{\Large\color{grayblue}Undefined-behavior prevention, ownership, validated inputs and review controls}\par\vspace{.4in}
\begin{keybox}\textbf{Standards basis:} ISO/IEC 9899:2024 (C23). Independent educational and engineering companion. Not an official ISO interpretation, full specification, certification or conformance assessment.\end{keybox}
\vfill{\large\bfseries C23 Software Engineering and Language Standards}\par\vspace{.15in}{\color{grayblue}Prepared October 2026 \quad | \quad Source edition: October 2024}\par
\end{titlepage}
\pagenumbering{roman}\tableofcontents\clearpage\pagenumbering{arabic}

\chapter*{Sources and interpretation key}
\addcontentsline{toc}{chapter}{Sources and interpretation key}
\textbf{Primary source:} ISO/IEC 9899:2024(en), \emph{Information technology --- Programming languages --- C}, fifth edition (2024). Clause and annex identifiers refer to the supplied licensed standard. This handbook contains original explanations, designs and examples; it does not reproduce the standard. The licensed standard remains the authoritative source for exact wording and conformance decisions.\par
\textbf{Secondary engineering sources:} \url{https://gcc.gnu.org/onlinedocs/gcc/Standards.html}; \url{https://clang.llvm.org/c_status.html}; \url{https://cmake.org/cmake/help/latest/prop_tgt/C_STANDARD.html}. Compiler support and diagnostics are release-specific: verify the installed compiler, target and library before relying on a feature.\par
\begin{keybox}\textbf{Interpretation categories.} \textbf{[STANDARD]} summarizes a requirement, constraint, semantic rule or defined facility and cites the applicable clause. \textbf{[OPTIONAL]} describes conditional or implementation-dependent support. \textbf{[GUIDANCE]} records an engineering recommendation, not an ISO mandate. \textbf{[TEST]} is an original check or exercise; a passing test is evidence but is not proof of full conformity.\end{keybox}
\textbf{Copyright boundary:} The supplied PDF is individually licensed. Do not redistribute it with this suite; distribute only these independently written companion materials in accordance with applicable rights.\par


\chapter{Scope and security decision framework}
ISO/IEC 9899:2024 defines language behavior, not a complete security assurance program. Security policies in this handbook are original engineering guidance informed by the specified behavior. CWE mappings are advisory vulnerability-class mappings, not source clauses or ISO requirements.\par
\begin{keybox}\textbf{Security invariant review:} check (1) the object exists, (2) lifetime remains valid, (3) bounds and alignment hold, (4) type conversions preserve intended values, (5) ownership is clear, (6) synchronization is sufficient, and (7) error outcomes are handled without inventing undefined behavior.\end{keybox}
\chapter{Security taxonomy and review prompts}
\small\begin{longtable}{p{.24\linewidth}p{.22\linewidth}p{.45\linewidth}}\toprule Hazard & Advisory CWE & Recommended control\\\midrule\endhead
\textbf{Numeric wrap and truncation} & CWE-190 / CWE-681 & Compute sizes with checked math; assert bounds before narrowing; test upper limits.\\\hline
\textbf{Out-of-bounds access} & CWE-125 / CWE-787 & Pass buffer extent and requested count; validate before pointer arithmetic.\\\hline
\textbf{Use after free} & CWE-416 & Define ownership and invalidation; avoid sharing unchecked borrowed pointers.\\\hline
\textbf{Double free} & CWE-415 & Enforce one owner and single destruction; nulling aliases is not a substitute.\\\hline
\textbf{Null dereference} & CWE-476 & Validate contract boundaries before dereferencing nullable inputs.\\\hline
\textbf{Uninitialized reads} & CWE-457 & Initialize local state and record all paths to first use.\\\hline
\textbf{Format string misuse} & CWE-134 & Keep format constant and pass untrusted strings through \%s with bounds.\\\hline
\textbf{Integer sign mismatch} & CWE-195 & Review signed-to-unsigned conversions around length and size checks.\\\hline
\textbf{Incorrect type punning} & CWE-843 & Use memcpy for byte copies; define endian handling and effective-type contract.\\\hline
\textbf{Misaligned access} & CWE-119 family & Use correctly aligned storage and avoid unaligned cast dereference.\\\hline
\textbf{Unexpected macro side effects} & CWE-628 family & Prefer functions over expression-repeating macros.\\\hline
\textbf{Improper locale/encoding handling} & CWE-176 / CWE-20 & Validate text encoding separately from byte length.\\\hline
\textbf{Races and stale state} & CWE-362 & Use atomic or mutex synchronization with explicit happens-before reasoning.\\\hline
\textbf{Invalid variadic arguments} & CWE-686 & Keep formats, promoted argument types and va\_arg extraction aligned.\\\hline
\textbf{Resource exhaustion} & CWE-400 & Cap allocations and IO sizes before calls; account for multiplicative growth.\\\hline
\textbf{Secret residual data} & CWE-226 & Minimize secret lifetime; use explicit erasure when available and appropriate.\\\hline
\textbf{Unexpected reachability} & CWE-758 family & Do not encode runtime fallible conditions as unreachable.\\\hline
\bottomrule\end{longtable}\normalsize
\chapter{Security review cards}
\section{Numeric wrap and truncation}
\textbf{Risk:} CWE-190 / CWE-681. \textbf{Standards basis:} \refcl{6.3.1, 6.5; 7.20}.\par\guidance\ Compute sizes with checked math; assert bounds before narrowing; test upper limits.\par\textbf{Verification:} Construct a passing boundary case, a failing boundary case and an explicit ownership or type invariant; record whether static analysis, compiler diagnostics or runtime instrumentation detected the failure.
\section{Out-of-bounds access}
\textbf{Risk:} CWE-125 / CWE-787. \textbf{Standards basis:} \refcl{6.5.3, 6.5.7; 7.26}.\par\guidance\ Pass buffer extent and requested count; validate before pointer arithmetic.\par\textbf{Verification:} Construct a passing boundary case, a failing boundary case and an explicit ownership or type invariant; record whether static analysis, compiler diagnostics or runtime instrumentation detected the failure.
\section{Use after free}
\textbf{Risk:} CWE-416. \textbf{Standards basis:} \refcl{6.2.4; 7.24.4}.\par\guidance\ Define ownership and invalidation; avoid sharing unchecked borrowed pointers.\par\textbf{Verification:} Construct a passing boundary case, a failing boundary case and an explicit ownership or type invariant; record whether static analysis, compiler diagnostics or runtime instrumentation detected the failure.
\section{Double free}
\textbf{Risk:} CWE-415. \textbf{Standards basis:} \refcl{7.24.4}.\par\guidance\ Enforce one owner and single destruction; nulling aliases is not a substitute.\par\textbf{Verification:} Construct a passing boundary case, a failing boundary case and an explicit ownership or type invariant; record whether static analysis, compiler diagnostics or runtime instrumentation detected the failure.
\section{Null dereference}
\textbf{Risk:} CWE-476. \textbf{Standards basis:} \refcl{6.5.4}.\par\guidance\ Validate contract boundaries before dereferencing nullable inputs.\par\textbf{Verification:} Construct a passing boundary case, a failing boundary case and an explicit ownership or type invariant; record whether static analysis, compiler diagnostics or runtime instrumentation detected the failure.
\section{Uninitialized reads}
\textbf{Risk:} CWE-457. \textbf{Standards basis:} \refcl{6.2.6; 6.7.11}.\par\guidance\ Initialize local state and record all paths to first use.\par\textbf{Verification:} Construct a passing boundary case, a failing boundary case and an explicit ownership or type invariant; record whether static analysis, compiler diagnostics or runtime instrumentation detected the failure.
\section{Format string misuse}
\textbf{Risk:} CWE-134. \textbf{Standards basis:} \refcl{7.23.6}.\par\guidance\ Keep format constant and pass untrusted strings through \%s with bounds.\par\textbf{Verification:} Construct a passing boundary case, a failing boundary case and an explicit ownership or type invariant; record whether static analysis, compiler diagnostics or runtime instrumentation detected the failure.
\section{Integer sign mismatch}
\textbf{Risk:} CWE-195. \textbf{Standards basis:} \refcl{6.3.1}.\par\guidance\ Review signed-to-unsigned conversions around length and size checks.\par\textbf{Verification:} Construct a passing boundary case, a failing boundary case and an explicit ownership or type invariant; record whether static analysis, compiler diagnostics or runtime instrumentation detected the failure.
\section{Incorrect type punning}
\textbf{Risk:} CWE-843. \textbf{Standards basis:} \refcl{6.2.6, 6.5}.\par\guidance\ Use memcpy for byte copies; define endian handling and effective-type contract.\par\textbf{Verification:} Construct a passing boundary case, a failing boundary case and an explicit ownership or type invariant; record whether static analysis, compiler diagnostics or runtime instrumentation detected the failure.
\section{Misaligned access}
\textbf{Risk:} CWE-119 family. \textbf{Standards basis:} \refcl{6.2.8}.\par\guidance\ Use correctly aligned storage and avoid unaligned cast dereference.\par\textbf{Verification:} Construct a passing boundary case, a failing boundary case and an explicit ownership or type invariant; record whether static analysis, compiler diagnostics or runtime instrumentation detected the failure.
\section{Unexpected macro side effects}
\textbf{Risk:} CWE-628 family. \textbf{Standards basis:} \refcl{6.10.5}.\par\guidance\ Prefer functions over expression-repeating macros.\par\textbf{Verification:} Construct a passing boundary case, a failing boundary case and an explicit ownership or type invariant; record whether static analysis, compiler diagnostics or runtime instrumentation detected the failure.
\section{Improper locale/encoding handling}
\textbf{Risk:} CWE-176 / CWE-20. \textbf{Standards basis:} \refcl{6.2.9; 7.11, 7.30}.\par\guidance\ Validate text encoding separately from byte length.\par\textbf{Verification:} Construct a passing boundary case, a failing boundary case and an explicit ownership or type invariant; record whether static analysis, compiler diagnostics or runtime instrumentation detected the failure.
\section{Races and stale state}
\textbf{Risk:} CWE-362. \textbf{Standards basis:} \refcl{5.1.2.4; 7.17--7.28}.\par\guidance\ Use atomic or mutex synchronization with explicit happens-before reasoning.\par\textbf{Verification:} Construct a passing boundary case, a failing boundary case and an explicit ownership or type invariant; record whether static analysis, compiler diagnostics or runtime instrumentation detected the failure.
\section{Invalid variadic arguments}
\textbf{Risk:} CWE-686. \textbf{Standards basis:} \refcl{7.16; 7.23.6}.\par\guidance\ Keep formats, promoted argument types and va\_arg extraction aligned.\par\textbf{Verification:} Construct a passing boundary case, a failing boundary case and an explicit ownership or type invariant; record whether static analysis, compiler diagnostics or runtime instrumentation detected the failure.
\section{Resource exhaustion}
\textbf{Risk:} CWE-400. \textbf{Standards basis:} \refcl{7.24}.\par\guidance\ Cap allocations and IO sizes before calls; account for multiplicative growth.\par\textbf{Verification:} Construct a passing boundary case, a failing boundary case and an explicit ownership or type invariant; record whether static analysis, compiler diagnostics or runtime instrumentation detected the failure.
\section{Secret residual data}
\textbf{Risk:} CWE-226. \textbf{Standards basis:} \refcl{7.26}.\par\guidance\ Minimize secret lifetime; use explicit erasure when available and appropriate.\par\textbf{Verification:} Construct a passing boundary case, a failing boundary case and an explicit ownership or type invariant; record whether static analysis, compiler diagnostics or runtime instrumentation detected the failure.
\section{Unexpected reachability}
\textbf{Risk:} CWE-758 family. \textbf{Standards basis:} \refcl{7.21.2}.\par\guidance\ Do not encode runtime fallible conditions as unreachable.\par\textbf{Verification:} Construct a passing boundary case, a failing boundary case and an explicit ownership or type invariant; record whether static analysis, compiler diagnostics or runtime instrumentation detected the failure.
\chapter{Secure programming patterns in C23}
\section{Length-aware copy API}
\begin{lstlisting}
#include <stdbool.h>
#include <stddef.h>
#include <string.h>
/* Preconditions: dst/src designate extents of at least cap/len. */
bool copy_bytes(void *dst, size_t cap,
                const void *src, size_t len) {
    if (len > cap) return false;
    if (len == 0) return true;
    if (dst == NULL || src == NULL) return false;
    /* Contract additionally requires nonoverlap. */
    memcpy(dst, src, len);
    return true;
}
\end{lstlisting}
\guidance\ This function does not dynamically prove object extents or overlap; it shifts responsibility to the API contract. Use \texttt{memmove} where overlap is permitted and specify expected behavior.
\section{Overflow-aware allocation}
\begin{lstlisting}
#include <stdint.h>
#include <stddef.h>
#include <stdlib.h>
void *safe_new_array(size_t n, size_t size) {
    if (size == 0 || n == 0 || n > SIZE_MAX / size) return NULL;
    return malloc(n * size);
}
\end{lstlisting}
\guidance\ A request that fits in \texttt{size\_t} may still exhaust resources; impose a domain-specific upper bound.
\section{Safe replacement of \texttt{realloc}}
\begin{lstlisting}
#include <stdlib.h>
#include <stdbool.h>
bool reserve(void **owned, size_t bytes) {
    if (!owned || bytes == 0) return false;
    void *next = realloc(*owned, bytes);
    if (!next) return false;
    *owned = next;
    return true;
}
\end{lstlisting}
\standard\ A failed nonzero allocation request must not cause the caller to lose its prior pointer. \refcl{7.24.4}.\guidance\ On success prior interior pointers may no longer be valid, even if the numerical address appears unchanged.
\section{Return-value and bounds contract for parsing}
\begin{lstlisting}
#include <errno.h>
#include <limits.h>
#include <stdbool.h>
#include <stdlib.h>
bool parse_nonnegative_int(const char *text, int *out) {
    if (text == NULL || out == NULL || *text == '\0') return false;
    errno = 0;
    char *end = NULL;
    long value = strtol(text, &end, 10);
    if (errno == ERANGE || end == text || *end != '\0' ||
        value < 0 || value > INT_MAX) return false;
    *out = (int)value;
    return true;
}
\end{lstlisting}
\guidance\ An application expecting only ASCII decimal digits may need additional lexical validation; \texttt{strtol} permits leading whitespace and an optional sign.
\section{Concurrency guard with a mutex}
\begin{lstlisting}
#include <threads.h>
struct Counter { mtx_t lock; unsigned value; };
int add_one(struct Counter *c) {
    if (mtx_lock(&c->lock) != thrd_success) return -1;
    c->value += 1u; /* unsigned wrap is defined; domain may forbid it */
    if (mtx_unlock(&c->lock) != thrd_success) return -1;
    return 0;
}
\end{lstlisting}
\guidance\ Initialization, shutdown and a maximum-value policy are required elsewhere; this fragment demonstrates the locking boundary, not a complete service.
\section{Explicit erasure and reachability}
\standard\ \texttt{memset\_explicit} appears in \refcl{7.26}; \texttt{unreachable} is specified in \refcl{7.21.2}.\guidance\ Use an explicit erasure facility only after target support has been established; it does not prevent secret copies from existing elsewhere. Never use \texttt{unreachable} where invalid input could actually arrive.
\chapter{Assessment workflow}
\begin{enumerate}
\item \textbf{Architecture pass:} enumerate trust boundaries, parsers, integer lengths and binary formats.
\item \textbf{Ownership pass:} build an allocation/free diagram with aliases and error paths.
\item \textbf{Type pass:} examine integer promotions, casts, bitfields and signed-to-unsigned boundaries.
\item \textbf{Memory pass:} prove extents, alignment, nullability, lifetime and overlap.
\item \textbf{Concurrency pass:} document shared mutable state and each synchronization relationship.
\item \textbf{Library pass:} check function argument domains, format strings, locale and resource limits.
\item \textbf{Verification pass:} combine warnings, SAST, negative tests, sanitizers and fuzzing.
\item \textbf{Evidence pass:} record findings, severity rationale, CWE, clause, fix, test and reviewer.
\end{enumerate}
\chapter{Secure C23 PR acceptance criteria}
\begin{itemize}
\item No new known reachable undefined behavior introduced; compile warnings reviewed and justified.
\item Every new pointer interface has documented extent and ownership; every allocation has checked sizing and cleanup.
\item Test cases cover minimum, maximum, zero, invalid, repeated and partial-failure paths.
\item No ignored critical findings; exceptional behavior requires a documented, approved disposition.
\item Dependencies and build flags are pinned and reproducible; optional C23 features are gated by capability tests.
\end{itemize}
\begin{warnbox}\textbf{Scope boundary.} A project can satisfy the ISO C language rules and still contain exploitable logic flaws or unsafe application behavior. Conversely, a safety coding policy can prohibit otherwise valid C23 constructs. Do not conflate these judgments.\end{warnbox}

\end{document}
